% Generated from validated synthesis. Do not hand-edit.
\tableofcontents
\medskip
\section*{Retrospective Signals}
% half-year final frontmatter compaction
\begingroup\small
\addcontentsline{toc}{section}{Retrospective Signals}
\vspace{0.08em}
\noindent\textbf{対話基盤がテキストの外へ広がった} GPT-4o、Gemini 1.5、Claude 3系を通して、長文脈・画像・音声・低遅延を統合したinteraction layerがフロンティア競争の中心になった。 \hfill{\footnotesize p.~\pageref{pkg:frontier-interaction}}\par
\vspace{0.08em}
\noindent\textbf{Openはweight公開からstackの再現性へ広がった} Llama 3、OLMo、DeepSeek-V2などにより、weightsだけでなくlicense、training artifacts、data、deploymentを分けて比較する必要が明確になった。 \hfill{\footnotesize p.~\pageref{pkg:open-ecosystem}}\par
\vspace{0.08em}
\noindent\textbf{配置制約と計算制約がmodel設計を変えた} Phi-3やApple Foundation Modelsの端末・プライベート環境への配置と、Jamba・Mamba-2・BitNet・LongRoPEのarchitecture探索が、parameter scale以外の設計軸を前景化した。 \hfill{\footnotesize p.~\pageref{pkg:edge-private} / p.~\pageref{pkg:architecture-efficiency}}\par
\vspace{0.08em}
\noindent\textbf{生成は行動とメディアへ分岐した} RAG、tool use、computer environmentへの接続と、Sora・Veo・Stable Diffusion 3・Klingに代表される生成メディアが、それぞれ独立した能力面として立ち上がった。 \hfill{\footnotesize p.~\pageref{pkg:action-rag-agents} / p.~\pageref{pkg:generative-media}}\par
\vspace{0.08em}
\noindent\textbf{能力を成立させるControl・Serving・Dataが技術史の一部になった} jailbreak、behavior specification、interpretability、NIM、FineWebが示したのは、強いmodelだけでなく制御・runtime・data pipelineまで含めて実用能力を考える必要性だった。 \hfill{\footnotesize p.~\pageref{pkg:safety-control} / p.~\pageref{pkg:systems-data} / p.~\pageref{pkg:conclusion}}\par
\begin{claimboundary}[Retrospective scope]
本号は2024年前期を後日確認可能になった一次情報も用いて再構成するRetrospective Specialである。Coverage Periodと制作時点を同一視せず、vendor / project / author claimの境界は各記事とTechnical Notesで明示する。
\end{claimboundary}
\endgroup
